\chapter{Implementation}

\section{Development Environment}
\begin{table}[H]\centering
\caption{Development Environment}\label{tab:devenv}
\begin{tabularx}{\textwidth}{|p{4.5cm}|X|}\hline
\textbf{Item} & \textbf{Details} \\ \hline
Operating system & Windows 11 (development); Linux containers (deployment) \\ \hline
Backend language and framework & Python 3.12+, Django 6.0, Django REST Framework 3.17, SimpleJWT 5.5 \\ \hline
Frontend & Node.js, Vue.js 3.5, Vue Router 4, Vite, Tailwind CSS 4, Reka UI, vue-pdf-embed \\ \hline
Database & PostgreSQL 16 with pgvector 0.8 (Docker image \code{pgvector/pgvector:pg16}) \\ \hline
Document processing & pypdfium2 5.7, python-docx 1.2, Tesseract OCR 5 via pytesseract \\ \hline
Machine learning & sentence-transformers 5.3, Hugging Face Transformers 5.5, PyTorch (CPU), scikit-learn, datasets \\ \hline
LLM providers & Groq (\code{openai/gpt-oss-20b}); Google Gemini (fallback and vision) \\ \hline
Editor and tools & Visual Studio Code, Git and GitHub, Docker Desktop, Google Colab (GPU training) \\ \hline
\end{tabularx}
\end{table}

\section{Core Module Implementation}
\subsection{Module 1: Organization and Access Management}
Registration (\code{POST /api/auth/register-organization/}) creates the organization, the user and an
Admin membership in one database transaction. Login uses SimpleJWT and writes a LOGIN audit event.
Permission classes (\code{IsOrgMember}, \code{IsOrgAdmin}, \code{CanUpload}, \code{CanVerify},
\code{CanViewAudit}) protect the API, and \code{visible\_documents(user)} implements the
document-level rule of FR-1.3:
\begin{lstlisting}[language=Python]
def visible_documents(user):
    m = getattr(user, 'membership', None)
    if m is None:
        return Document.objects.none()
    qs = Document.objects.filter(organization_id=m.organization_id, is_archived=False)
    if m.sees_everything:                        # Admin, Auditor
        return qs
    rule = Q(access_level='ORG') | Q(uploaded_by=user) | Q(allowed_users=user)
    if m.department_id:
        rule |= Q(access_level='DEPARTMENT', department_id=m.department_id)
    return Document.objects.filter(id__in=qs.filter(rule).values('id'))
\end{lstlisting}

\subsection{Module 2: Document Ingestion and Processing}
The upload endpoint validates the file extension and size, stores the file as a new
\code{DocumentVersion}, computes its SHA-256 hash and creates a \code{ProcessingJob}. The
\code{run\_worker} command claims jobs and calls \code{process\_version}, which runs five logged steps:
\begin{enumerate}
\item \textbf{Extract.} PDFs are opened with pypdfium2. A page with a usable text layer is taken as is;
otherwise it is rendered at 300\,dpi and passed to OCR. DOCX files are read with python-docx; images go
straight to OCR. The OCR step runs Tesseract and computes the mean word confidence; if the document is
marked handwritten, the confidence is below 70, or almost no words are found, the page image is sent to
the vision-language model with an instruction to transcribe exactly and mark unreadable words.
\item \textbf{Classify.} The first pages are sent to the LLM, which returns the kind, document type and a
summary in JSON. A keyword heuristic is used if the LLM is unavailable; it recognizes statutes (``Act'',
``Ordinance'', ``Code'' followed by a year) as reference material.
\item \textbf{Chunk and embed.} The structure-aware chunker (Algorithm~\ref{alg:chunk}) produces chunks
with section paths and page ranges; each is embedded with all-MiniLM-L6-v2 together with its section
path, and a weighted \code{tsvector} is computed in SQL.
\item \textbf{Label clauses.} For records, the fine-tuned classifier labels each chunk if its confidence
is at least 0.5.
\item \textbf{Extract fields.} For records, the LLM receives the paged text (up to 22,000 characters; for
longer documents the beginning and end are kept, because signatures and dates are usually there) and
returns the eleven fields with value, page, quote and confidence. Each quote is grounded
(Algorithm~\ref{alg:ground}).
\end{enumerate}
Failures are retried up to three times. When a version becomes READY it becomes the current version and
the uploader is notified.

\subsection{Module 3: Records Register and Verification}
The records endpoint returns one row per visible RECORD document with its key fields from the current
version. The review queue returns NEEDS\_REVIEW fields of visible documents. Verifying or correcting a
field updates its status (and value), and creates a \code{FieldEdit} row and an audit event in the same
transaction. Users can also add a field the model missed.

\subsection{Module 4: Legal Assistant}
\code{AskView} (throttled to 20 requests per minute) calls \code{engine.answer(user, text,
document\_ids)}. The router returns a plan; the engine runs the matching path:
\begin{itemize}
\item \textbf{LOOKUP} filters \code{RecordField} rows of the current versions of visible documents by
document type (using a table of aliases, e.g.\ NDA $\rightarrow$ ``non-disclosure'', ``confidentiality
agreement''), party name and field; date ranges use \code{value\_date}. The answer is composed from a
template (``The NDA \emph{X} was signed on \emph{12 March 2025} (page 3)''), so no LLM is needed and
the answer takes about a second.
\item \textbf{KNOWLEDGE} retrieves the eight best chunks with hybrid search (Algorithm~\ref{alg:rrf})
and asks the LLM to answer only from the numbered sources and cite them as [n].
\item \textbf{SEARCH} groups retrieved chunks by document, boosts chunks whose clause label matches the
query, and returns the documents with their best passages.
\end{itemize}
Every LLM call has a deadline. If the main provider fails or reports that its daily quota is used up,
the gateway switches to the fallback provider and skips the main one until the quota resets; if both
fail, the user still receives the most relevant passages. The query, plan, sources and latency are
saved to the conversation.

\subsection{Module 5: History, Audit and Notifications}
\code{audit.service.log()} is called from every view that changes or reveals data. The
\code{AuditEvent} model overrides \code{save()} and \code{delete()} to forbid changes after creation.
\code{notifications.service.run\_reminders()} runs hourly in the worker: it computes expiry dates and
notice deadlines (expiry minus the parsed notice period), finds those 90, 30, 7 or 0 days ahead, selects
the people who can see the document (uploader, department managers, admins) and creates notifications
with a unique \emph{dedupe key}, so each reminder is sent once.

\subsection{Module 6: Model Training and Evaluation}
\code{ml/train\_clause\_classifier.py} fine-tunes the clause classifier (Section~\ref{sec:finetune});
\code{ml/baselines.py} trains the TF-IDF baseline and runs the zero-shot LLM baseline; the
\code{evaluation/} package contains one script per component and \code{evaluation.report}, which writes
all results into \code{docs/EVALUATION.md}.

\section{Algorithm Implementation}

\begin{algorithm}[H]
\caption{Structure-aware chunking}\label{alg:chunk}
\begin{algorithmic}[1]
\Require pages $P = [(n, \mathit{text})]$, $\mathit{MAX}=1500$, $\mathit{MIN}=200$
\State $\mathit{stack} \gets [\,]$; $\mathit{sections} \gets [\,]$; $\mathit{cur} \gets$ new section
\For{each line $\ell$ on page $n$ in $P$}
  \If{$\ell$ matches a heading rule of level $k$ (PART/CHAPTER/ARTICLE, Section $n$, $n.n$, ``$n$.'' title)}
    \State append $\mathit{cur}$ to $\mathit{sections}$
    \State pop from $\mathit{stack}$ every heading with level $\geq k$; push $(k,\ell)$
    \State $\mathit{cur} \gets$ new section with path $=$ join($\mathit{stack}$, `` $>$ '')
  \Else
    \State add $(n,\ell)$ to $\mathit{cur}$
  \EndIf
\EndFor
\State append $\mathit{cur}$; merge sections shorter than $\mathit{MIN}$ into the previous one
\For{each section $s$}
  \State split $s$ on paragraphs, then sentences, into parts of at most $\mathit{MAX}$ characters with overlap
  \State emit chunks with $s$.path, $s$.heading and the first and last page of their lines
\EndFor
\end{algorithmic}
\end{algorithm}

\begin{algorithm}[H]
\caption{Permission-filtered hybrid retrieval with Reciprocal Rank Fusion}\label{alg:rrf}
\begin{algorithmic}[1]
\Require user $u$, query $q$, optional document ids $D$, $k=8$, $C=40$, $K_{\mathrm{RRF}}=60$
\State $V \gets$ ids of current versions of \Call{VisibleDocuments}{$u$} \Comment{access rule first}
\If{$D$ given} $V \gets V \cap$ versions of $D$ \EndIf \Comment{forbidden ids are dropped}
\State $e \gets$ \Call{Embed}{$q$}
\State $L_d \gets$ top $C$ chunks in $V$ by cosine distance to $e$ (HNSW, \code{ef\_search}=100)
\State $L_s \gets$ top $C$ chunks in $V$ matching any content word of $q$, ordered by \code{ts\_rank}
\For{each chunk $c$}
  \State $\mathit{score}(c) \gets \sum_{L \in \{L_d, L_s\}} \frac{1}{K_{\mathrm{RRF}} + \mathrm{rank}_L(c)}$
\EndFor
\State \Return the $k$ chunks with the highest score
\end{algorithmic}
\end{algorithm}

\begin{algorithm}[H]
\caption{Field grounding and review decision}\label{alg:ground}
\begin{algorithmic}[1]
\Require field $f$ (value, page, quote, confidence $c$), page texts $T$, threshold $\tau = 0.75$
\State $g \gets$ \Call{Ground}{$f$.quote, $T[f.\mathit{page}]$}
\If{not $g$} \Comment{the model may have cited the wrong page}
  \For{each page $p$} \If{\Call{Ground}{$f$.quote, $T[p]$}} $g \gets$ true; $f.\mathit{page} \gets p$; \textbf{break} \EndIf \EndFor
\EndIf
\If{not $g$} $c \gets c \times 0.5$ \EndIf
\If{$f$ is a date and cannot be parsed} $c \gets c \times 0.7$ \EndIf
\State $f.\mathit{status} \gets$ NEEDS\_REVIEW \textbf{if} $c < \tau$ or not $g$ or page source $\in$ \{OCR, HANDWRITING\} \textbf{else} AUTO
\Statex
\Function{Ground}{$q$, $t$} \Comment{normalize case, punctuation, spaces}
  \State \Return $q \subseteq t$ as a string, or ($|q| \geq 3$ words and $\geq 80\%$ of the words of $q$ occur in $t$)
\EndFunction
\end{algorithmic}
\end{algorithm}

\begin{algorithm}[H]
\caption{Query routing}\label{alg:route}
\begin{algorithmic}[1]
\Require question $q$, last three questions $H$ of the conversation
\State ask the LLM (JSON mode, deadline 12\,s) for \{intent, doc\_type, party, field, date\_from, date\_to, keywords\}, using $H$ to resolve ``it'' / ``that contract''
\If{the call fails or the intent is invalid}
  \State field $\gets$ first synonym rule that matches $q$ (e.g.\ ``signed'' $\rightarrow$ signed\_date)
  \State doc\_type $\gets$ record word in $q$ (NDA, lease, contract \dots); party $\gets$ capitalized name after ``with/between/for''
  \If{field and (doc\_type or party)} intent $\gets$ LOOKUP
  \ElsIf{$q$ starts with find/show/open/list/search} intent $\gets$ SEARCH
  \Else\ intent $\gets$ KNOWLEDGE \EndIf
\EndIf
\State replace string values such as ``null'' or ``N/A'' by none; validate the field key and ISO dates
\State \Return plan
\end{algorithmic}
\end{algorithm}

\section{Fine-tuning the Clause Classifier}\label{sec:finetune}
The supervisor asked for a fine-tuned model. Clause classification was chosen because it has a large
public labelled dataset and a clear use in the system (labels on the document page and a search boost).
\begin{itemize}
\item \textbf{Dataset.} LEDGAR~\cite{tuggener2020ledgar} from LexGLUE~\cite{chalkidis2022lexglue}:
contract provisions from SEC filings labelled with 100 clause types (60,000 training, 10,000
validation and 10,000 test provisions).
\item \textbf{Model.} \code{sentence-transformers/all-MiniLM-L6-v2}~\cite{wang2020minilm} (22 million
parameters, 6 layers) with a classification head, chosen because it runs quickly on a CPU at inference
time and shares its encoder family with the retrieval embeddings.
\item \textbf{Training.} Hugging Face \code{Trainer}, learning rate $5\times10^{-5}$, batch size 32,
maximum length 256 tokens, 3 epochs, warm-up ratio 0.06, weight decay 0.01, seed 42. Because a laptop
CPU was used, the reported run used 10,000 training provisions (36.7 minutes). A Google Colab notebook
(\code{ml/colab\_train\_clause\_classifier.ipynb}) trains on all 60,000 provisions with a GPU.
\item \textbf{Engineering note.} With Transformers 5.5, reloading intermediate checkpoints silently reset
some LayerNorm weights. Mid-run checkpoints were therefore disabled, the model is saved once at the end,
and the test set is evaluated on the \emph{reloaded} model so that the reported numbers are those of the
model the application actually uses.
\item \textbf{Baselines.} TF-IDF with logistic regression (10,000 and 60,000 training provisions) and a
zero-shot LLM that picks a label from the list.
\end{itemize}
Results are reported in Section~\ref{sec:eval-clauses}.

\section{External APIs / SDKs}
\begin{table}[H]\centering
\caption{External APIs and Third-Party Components Used}\label{tab:apis}
\begin{tabularx}{\textwidth}{|p{4.2cm}|X|p{3.6cm}|}\hline
\textbf{Component} & \textbf{Use in the system} & \textbf{Where} \\ \hline
Groq API (\code{gpt-oss-20b}) & Routing plans, classification, field extraction, written answers & \code{ai/llm.py} \\ \hline
Google Gemini API & Fallback text model; vision model for handwriting & \code{ai/llm.py}, \code{ai/ocr.py} \\ \hline
Tesseract OCR & Printed text recognition with word confidences & \code{ai/ocr.py} \\ \hline
sentence-transformers & Local 384-dimensional embeddings & \code{ai/embeddings.py} \\ \hline
Hugging Face Transformers & Clause classifier training and inference & \code{ml/}, \code{ai/clauses.py} \\ \hline
pgvector & Vector column, HNSW index, cosine distance in SQL & \code{documents/models.py} \\ \hline
pypdfium2 / python-docx & PDF text and rendering / DOCX reading & \code{pipeline/extract.py} \\ \hline
vue-pdf-embed, mammoth & PDF rendering with text layer; DOCX preview in the browser & frontend \\ \hline
marked, DOMPurify & Rendering answers as safe Markdown & frontend \\ \hline
\end{tabularx}
\end{table}
All LLM providers are called through their OpenAI-compatible Chat Completions endpoint, so switching
provider (including a local Ollama model for fully offline use) only requires changing
\code{LLM\_PROVIDER} and \code{LLM\_MODEL} in the environment file.

\section{User Interface Implementation}
\subsection{UI Design Approach}
The interface was redesigned from scratch after the first version was judged generic. It uses a calm,
light ``glass on paper'' style: translucent panels over a subtle paper-grain background, a single neutral
accent, and no colour gradients, so that documents and evidence remain the focus. Key decisions:
\begin{itemize}
\item \textbf{Evidence everywhere.} Every extracted value has a ``show evidence'' action that opens the
page with the quote highlighted; every answer lists numbered sources that do the same.
\item \textbf{Question box at the bottom.} The Assistant uses a docked input, like a chat, and shows the
steps it is taking (understanding the question, checking access, reading pages, writing the answer).
\item \textbf{Keyboard support.} A command palette (Ctrl+K) searches documents and actions; ``/''
focuses the question box; focus rings are visible.
\item \textbf{Role-aware navigation.} Review, Team and Audit appear only for roles that may use them.
\end{itemize}

\subsection{Screen-wise Implementation}
\subsubsection{Screen 1: Landing Page}
Introduces the product with a self-playing demonstration of a question being answered with a citation.
\screenshot{landing.jpg}{Landing Page}{fig:scr-landing}

\subsubsection{Screen 2: Sign-up and Login}
Sign-up registers a new organization and its administrator; login signs in existing members.
\screenshot{login.jpg}{Login Page}{fig:scr-login}

\subsubsection{Screen 3: Assistant}
The main screen. Users ask questions, optionally limited to chosen documents. Answers show the route
taken, the time, the written answer with numbered citations, and source cards.
\screenshot{assistant.jpg}{Assistant -- Answer with Citations}{fig:scr-assistant}

\subsubsection{Screen 4: Documents}
Lists visible documents with kind, type, status, access level and page count, with search, filters and
upload.
\screenshot{documents.jpg}{Documents Screen}{fig:scr-documents}

\subsubsection{Screen 5: Upload Modal}
Drag-and-drop upload of several files with title, department, access level, kind and the
``handwritten'' flag.
\screenshot{upload.jpg}{Upload Documents Modal}{fig:scr-upload}

\subsubsection{Screen 6: Document Detail}
Shows the document in an embedded viewer next to tabs for extracted facts (with evidence), clauses,
versions and comparison, extracted text and activity.
\screenshot{document_detail.jpg}{Document Detail: Extracted Facts with the Signing-date Evidence Highlighted}{fig:scr-detail}

Opening a source from an answer shows the cited passage of the reference document highlighted on its page
(Figure~\ref{fig:scr-source}).
\screenshot{source_highlight.jpg}{Source of an Answer Opened in the Contract Act with the Passage Highlighted}{fig:scr-source}

\subsubsection{Screen 7: Records Register}
\screenshot{records.jpg}{Records Register}{fig:scr-records}

\subsubsection{Screen 8: Review Queue}
\screenshot{review.jpg}{Review Queue}{fig:scr-review}

\subsubsection{Screen 9: Team Management}
\screenshot{team.jpg}{Team Management Screen}{fig:scr-team}

\subsubsection{Screen 10: Audit Log}
\screenshot{audit.jpg}{Audit Log Screen}{fig:scr-audit}

\subsubsection{Screen 11: Profile Settings}
\screenshot{profile.jpg}{Profile Settings}{fig:scr-profile}

\subsection{Navigation Flow Diagram}
\begin{figure}[H]\centering
\begin{adjustbox}{max width=\textwidth}
\begin{tikzpicture}[node distance=7mm and 9mm]
\node[box] (land) {Landing};
\node[box,right=of land] (login) {Login};
\node[box,below=of login] (reg) {Sign-up};
\node[box,right=14mm of login,fill=gray!8] (ask) {Assistant};
\node[box,right=of ask] (docs) {Documents};
\node[box,right=of docs] (det) {Document detail};
\node[box,below=of docs] (rec) {Records};
\node[box,below=of rec] (rev) {Review$^{\ast}$};
\node[box,above=of docs] (team) {Team$^{\ast}$};
\node[box,above=of ask] (aud) {Audit$^{\ast}$};
\node[box,below=of ask] (prof) {Profile};
\draw[flow] (land) -- (login);
\draw[flow] (land) |- (reg);
\draw[flow] (login) -- (ask);
\draw[flow] (reg) -| ($(login.east)!0.5!(ask.west)$) -- (ask);
\draw[flow] (ask) -- (docs);
\draw[flow] (docs) -- (det);
\draw[flow] (ask.north east) to[bend left=20] node[lbl,above]{source} (det.north west);
\draw[flow] (rec) -| (det);
\draw[flow] (rev) -| (det.south east);
\draw[flow] (ask) -- (aud);
\draw[flow] (ask) -- (prof);
\draw[flow] (ask) -- (team);
\draw[flow] (ask) -- (rec);
\end{tikzpicture}
\end{adjustbox}
\caption{Application Navigation Flow ($^{\ast}$ shown only to permitted roles)}\label{fig:navflow}
\end{figure}

\subsection{Responsiveness and Performance}
The frontend is built with Vite and loads each screen as a separate, lazily loaded bundle, so the first
page opens quickly. Heavy work (OCR, embeddings, model calls) never runs in the browser or in the
request that uploads a file; the interface polls the document status and updates when processing
finishes. The layout uses responsive breakpoints for laptop and tablet widths.

\section{Deployment}
\begin{table}[H]\centering
\caption{Deployment Details}\label{tab:deploy}
\begin{tabularx}{\textwidth}{|p{3cm}|X|}\hline
\textbf{Service} & \textbf{Configuration (docker-compose.yml)} \\ \hline
db & \code{pgvector/pgvector:pg16}, persistent volume, health check \\ \hline
backend & Python 3.12-slim image with Tesseract and CPU PyTorch; runs migrations and Gunicorn \\ \hline
worker & Same image; runs \code{python manage.py run\_worker} (processing and hourly reminders) \\ \hline
web & Node build of the Vue app served by Nginx, which also proxies \code{/api/} to the backend \\ \hline
Configuration & All secrets and provider settings in a \code{.env} file (never committed); \code{.env.example} documents every variable \\ \hline
\end{tabularx}
\end{table}
For development, the database runs in Docker while the API, worker and Vite dev server run locally.
A demonstration organization with users for every role and sample documents is created with
\code{python manage.py seed\_demo}.

\section{Chapter Summary}
This chapter described the development environment, the implementation of the six modules, the main
algorithms, the fine-tuning of the clause classifier, the external components, the user interface and
the deployment configuration. The next chapter tests and evaluates the system.
